% ==============================================================================
% CHAPITRE 2 : PUISSANCES DE 10 & ÉCRITURES SCIENTIFIQUES (VERSION ÉLÈVE)
% ==============================================================================
\chapter{Puissances de 10, Calculs \& Écritures Scientifiques}

% ------------------------------------------------------------------------------
% 1. ACTIVITÉ D'INVESTIGATION (SITUATION PROFESSIONNELLE / SCIENCES)
% ------------------------------------------------------------------------------
\begin{activite}{De l'infiniment petit à l'infiniment grand (Microélectronique \& Télécoms)}
Dans le secteur de l'informatique et des télécommunications, les ingénieurs manipulent des grandeurs extrêmes :
\begin{itemize}
    \item L'épaisseur d'un transistor dans un microprocesseur moderne est de \textbf{0{,}000 000 003 m} ($3\text{ nanomètres}$).
    \item La capacité totale des centres de données mondiaux dépasse \textbf{120 000 000 000 000 000 000 octets} ($120\text{ exaoctets}$).
    \item La vitesse de transmission d'un signal laser dans la fibre optique est d'environ \textbf{300 000 000 m/s}.
\end{itemize}

\begin{center}
\begin{tikzpicture}
    \node[draw=colPrimary, fill=slate50, rounded corners=6pt, inner sep=8pt, text width=\dimexpr\linewidth-28pt\relax] {
        \sffamily
        \textbf{Problématique :}
        \begin{enumerate}
            \item Pourquoi l'écriture décimale avec tous ces zéros est-elle difficile à utiliser lors des calculs ?
            \item Comment peut-on écrire ces nombres de manière plus compacte et lisible à l'aide des puissances de 10 ?
        \end{enumerate}
    };
\end{tikzpicture}
\end{center}

\vspace{1mm}
\noindent\textbf{Espace de recherche :}
\lignesreponse[4]
\end{activite}

% ------------------------------------------------------------------------------
% 2. COURS : PUISSANCES D'UN NOMBRE & PUISSANCES DE 10
% ------------------------------------------------------------------------------
\section{Puissances d'un Nombre Relatif \& Puissances de 10}

\begin{cadredef}[Puissance d'un nombre]
Soit $a$ un nombre relatif et $n$ un nombre entier positif ($n \geq 1$).
\begin{itemize}
    \item $a^n = \underbrace{a \times a \times \dots \times a}_{n \text{ facteurs}}$ \quad ($a^n$ se lit « $a$ puissance $n$ » ou « $a$ exposant $n$ »).
    \item Par convention : pour tout nombre $a \neq 0$, \textbf{$a^0 = 1$} \quad et \quad \textbf{$a^1 = a$}.
    \item Pour tout nombre $a \neq 0$ : \textbf{$a^{-n} = \dfrac{1}{a^n}$} \quad ($a^{-n}$ est l'\textbf{inverse} de $a^n$).
\end{itemize}
\end{cadredef}

\begin{cadreexemple}[Exemples à compléter ensemble en classe]
\begin{itemize}
    \item $2^4 = 2 \times 2 \times 2 \times 2 = \pointille[2cm]$ \hfill $(-3)^2 = (-3) \times (-3) = \pointille[2cm]$
    \item $5^0 = \pointille[1.5cm]$ \hfill $5^{-1} = \dfrac{1}{5} = \pointille[1.5cm]$ \hfill $2^{-3} = \dfrac{1}{2^3} = \dfrac{1}{\pointille[1cm]} = \pointille[1.5cm]$
\end{itemize}
\end{cadreexemple}

\begin{cadreprop}[Puissances de 10 positives et négatives]
Soit $n$ un nombre entier positif :
\begin{itemize}
    \item \textbf{Exposant positif} : $10^n = 1\underbrace{00\dots0}_{n\text{ zéros}}$ \quad (le chiffre 1 suivi de $n$ zéros).
    \item \textbf{Exposant négatif} : $10^{-n} = \dfrac{1}{10^n} = \underbrace{0{,}00\dots0}_{n\text{ zéros}}1$ \quad ($n$ zéros au total, virgule après le 1\textsuperscript{er} zéro).
\end{itemize}
\end{cadreprop}

\begin{cadremethode}[Écriture décimale des puissances de 10]
Compléter le tableau de correspondance ci-dessous :
\begin{center}
\renewcommand{\arraystretch}{1.3}
\begin{tabularx}{\linewidth}{|c|X|c|X|}
\hline
\rowcolor{slate800} \textcolor{white}{\textbf{Puissance}} & \textcolor{white}{\textbf{Écriture décimale}} & \textcolor{white}{\textbf{Puissance}} & \textcolor{white}{\textbf{Écriture décimale}} \\
\hline
$10^1$ & \textbf{10} & $10^{-1}$ & \textbf{0,1} \\
\hline
$10^2$ & & $10^{-2}$ & \\
\hline
$10^3$ & \textbf{1\,000} (mille) & $10^{-3}$ & \textbf{0,001} (millième) \\
\hline
$10^6$ & & $10^{-6}$ & \\
\hline
$10^9$ & \textbf{1\,000\,000\,000} (milliard) & $10^{-9}$ & \\
\hline
\end{tabularx}
\end{center}
\end{cadremethode}

\begin{cadreprop}[Préfixes du Système International (Unités usuelles)]
\begin{center}
\renewcommand{\arraystretch}{1.2}
\small
\begin{tabular}{|c|c|c||c|c|c|}
\hline
\rowcolor{slate200} \textbf{Préfixe} & \textbf{Symbole} & \textbf{Puissance} & \textbf{Préfixe} & \textbf{Symbole} & \textbf{Puissance} \\
\hline
téra & T & $10^{12}$ & déci & d & $10^{-1}$ \\
\hline
giga & G & $10^9$ & centi & c & $10^{-2}$ \\
\hline
méga & M & $10^6$ & milli & m & $10^{-3}$ \\
\hline
kilo & k & $10^3$ & micro & $\mu$ & $10^{-6}$ \\
\hline
hecto & h & $10^2$ & nano & n & $10^{-9}$ \\
\hline
\end{tabular}
\end{center}
\end{cadreprop}

% ------------------------------------------------------------------------------
% 3. COURS : RÈGLES DE CALCUL SUR LES PUISSANCES DE 10
% ------------------------------------------------------------------------------
\section{Règles de Calcul sur les Puissances de 10}

\begin{cadreprop}[Formules fondamentales sur les puissances de 10]
Pour tous nombres entiers relatifs $m$ et $n$ :
\begin{center}
\renewcommand{\arraystretch}{1.8}
\begin{tabularx}{\linewidth}{X|X|X}
\rowcolor{colDefBg}
\centering \textbf{Produit} & \centering \textbf{Quotient} & \centering \textbf{Puissance de puissance} \tabularnewline
\centering $10^m \times 10^n = \mathbf{10^{m+n}}$ & \centering $\dfrac{10^m}{10^n} = \mathbf{10^{m-n}}$ & \centering $\left(10^m\right)^n = \mathbf{10^{m \times n}}$ \tabularnewline
\end{tabularx}
\end{center}
\end{cadreprop}

\begin{cadremethode}[Application des règles de calcul]
Compléter les calculs modèles suivants :
\begin{itemize}
    \item $A = 10^4 \times 10^3 = 10^{4+3} = \mathbf{10^7}$ \hfill $B = 10^7 \times 10^{-2} = 10^{\pointille[1.5cm]} = \pointille[1.5cm]$
    \item $C = \dfrac{10^8}{10^5} = 10^{8-5} = \mathbf{10^3}$ \hfill $D = \dfrac{10^3}{10^{-4}} = 10^{3 - (-4)} = 10^{\pointille[1.5cm]} = \pointille[1.5cm]$
    \item $E = \left(10^3\right)^4 = 10^{3 \times 4} = \mathbf{10^{12}}$ \hfill $F = \left(10^{-2}\right)^5 = 10^{\pointille[2cm]} = \pointille[1.5cm]$
\end{itemize}
\end{cadremethode}

% ------------------------------------------------------------------------------
% 4. COURS : ÉCRITURE SCIENTIFIQUE
% ------------------------------------------------------------------------------
\section{Notation Scientifique \& Ordres de Grandeur}

\begin{cadredef}[Écriture Scientifique]
L'\textbf{écriture scientifique} (ou notation scientifique) d'un nombre décimal non nul est son unique écriture de la forme :
\[ a \times 10^n \]
où :
\begin{itemize}
    \item $a$ est un nombre décimal ayant \textbf{un seul chiffre non nul avant la virgule} ($1 \leq a < 10$ ou $-10 < a \leq -1$).
    \item $n$ est un \textbf{nombre entier relatif} (positif ou négatif).
\end{itemize}
\end{cadredef}

\begin{cadremethode}[Méthode pas à pas : Mettre un nombre en écriture scientifique]
\begin{enumerate}
    \item On place la virgule juste après le premier chiffre non nul pour obtenir le nombre $a$.
    \item On compte le nombre de rangs dont la virgule s'est déplacée pour trouver l'exposant $n$ :
    \begin{itemize}
        \item Déplacement vers la \textbf{gauche} $\implies$ exposant \textbf{positif} ($+n$).
        \item Déplacement vers la \textbf{droite} $\implies$ exposant \textbf{négatif} ($-n$).
    \end{itemize}
\end{enumerate}

\textbf{Exemples à compléter :}
\begin{itemize}
    \item $45\,000 = 4{,}5 \times 10^4$ \hfill $728\,000\,000 = \pointille[2cm] \times 10^{\pointille[0.8cm]}$
    \item $0{,}003\,8 = 3{,}8 \times 10^{-3}$ \hfill $0{,}000\,005\,12 = \pointille[2cm] \times 10^{\pointille[0.8cm]}$
\end{itemize}
\end{cadremethode}

\begin{cadreretenu}[L'Essentiel du Chapitre 2]
\begin{itemize}
    \item $10^n = 100\dots0$ ($n$ zéros) \quad et \quad $10^{-n} = 0{,}00\dots01$ ($n$ zéros au total).
    \item $10^0 = 1$ \quad et \quad $10^1 = 10$.
    \item Règles de calcul : $10^m \times 10^n = 10^{m+n}$ \quad ; \quad $\dfrac{10^m}{10^n} = 10^{m-n}$ \quad ; \quad $(10^m)^n = 10^{m \times n}$.
    \item Écriture scientifique : $a \times 10^n$ avec $1 \leq a < 10$.
    \item Savoir convertir les unités avec les préfixes : $\text{kilo } (10^3)$, $\text{méga } (10^6)$, $\text{giga } (10^9)$, $\text{milli } (10^{-3})$, $\text{micro } (10^{-6})$, $\text{nano } (10^{-9})$.
\end{itemize}
\end{cadreretenu}

% ------------------------------------------------------------------------------
% 5. QUESTIONS FLASH / AUTOMATISMES
% ------------------------------------------------------------------------------
\section{Questions Flash / Automatismes}

\begin{automatisme}[8 Questions Flash d'Échauffement]
\begin{enumerate}[label=\textbf{Q\arabic*.}]
    \item Donner l'écriture décimale de $10^5$ : \pointille[5cm]
    \item Donner l'écriture décimale de $10^{-3}$ : \pointille[5cm]
    \item Écrire sous la forme d'une seule puissance de 10 : $10^4 \times 10^5 = \pointille[3cm]$
    \item Écrire sous la forme d'une seule puissance de 10 : $\dfrac{10^9}{10^2} = \pointille[3cm]$
    \item Écrire sous la forme d'une seule puissance de 10 : $\left(10^3\right)^2 = \pointille[3cm]$
    \item Donner l'écriture scientifique de $85\,000$ : \pointille[5cm]
    \item Donner l'écriture scientifique de $0{,}000\,72$ : \pointille[5cm]
    \item Convertir $4{,}5\text{ kW}$ en Watts sous forme de puissance de 10 : $4{,}5\text{ kW} = \pointille[3cm]\text{ W}$
\end{enumerate}
\end{automatisme}

% ------------------------------------------------------------------------------
% 6. TRAVAUX DIRIGÉS (TD) - EXERCICES D'ENTRAÎNEMENT PROGRESSIFS
% ------------------------------------------------------------------------------
\section{Travaux Dirigés (TD) : Exercices d'Entraînement}

\begin{exercice}[2 pts]{\capSApproprier}{Écritures décimales et puissances de 10}
Écrire chaque nombre sous forme décimale :
\begin{enumerate}
    \item $A = 10^6 = \pointille[6cm]$
    \item $B = 10^{-4} = \pointille[6cm]$
    \item $C = 3{,}4 \times 10^3 = \pointille[6cm]$
    \item $D = 8{,}15 \times 10^{-2} = \pointille[6cm]$
\end{enumerate}
\end{exercice}

\begin{exercice}[3 pts]{\capRealiser}{Calculs avec les puissances de 10}
Écrire les expressions suivantes sous la forme $10^p$, où $p$ est un nombre entier relatif :
\begin{enumerate}
    \item $E = 10^8 \times 10^4 = \pointille[6cm]$
    \item $F = 10^5 \times 10^{-7} = \pointille[6cm]$
    \item $G = \dfrac{10^{11}}{10^4} = \pointille[6cm]$
    \item $H = \dfrac{10^2}{10^{-5}} = \pointille[6cm]$
    \item $I = \left(10^4\right)^{-3} = \pointille[6cm]$
    \item $J = \dfrac{10^5 \times 10^3}{10^2} = \pointille[6cm]$
\end{enumerate}
\end{exercice}

\begin{exercice}[3 pts]{\capRealiser}{Passage en écriture scientifique}
Donner l'écriture scientifique de chacun des nombres suivants :
\begin{enumerate}
    \item $K = 540\,000 = \pointille[6cm]$
    \item $L = 0{,}000\,082 = \pointille[6cm]$
    \item $M = 37{,}5 \times 10^4 = \pointille[6cm]$
    \item $N = 0{,}004\,6 \times 10^{-3} = \pointille[6cm]$
\end{enumerate}
\end{exercice}

\begin{exercice}[4 pts]{\capAnalyser}{Conversions et préfixes (Mécanique \& Électronique)}
Un technicien en usinage de précision mesure les cotes de différentes pièces mécaniques :
\begin{itemize}
    \item Pièce A : épaisseur de $25\text{ mm}$.
    \item Pièce B : jeu mécanique de $40\text{ }\mu\text{m}$ (micromètres).
    \item Pièce C : rugosité de surface de $150\text{ nm}$ (nanomètres).
\end{itemize}

\begin{enumerate}
    \item Exprimer chacune de ces dimensions en \textbf{mètres} sous forme d'une \textbf{écriture scientifique} :
    \begin{itemize}
        \item Épaisseur Pièce A : \pointille[6cm]
        \item Jeu Pièce B : \pointille[6cm]
        \item Rugosité Pièce C : \pointille[6cm]
    \end{itemize}
    \item Ranger ces trois dimensions par ordre croissant (de la plus petite à la plus grande).
\end{enumerate}
\lignesreponse[4]
\end{exercice}

\begin{exercice}[4 pts]{\capAnalyser}{Stockage Numérique \& Téléchargement (Informatique)}
Un disque dur externe SSD possède une capacité de stockage de \textbf{2 To} (téraoctets).\\
On rappelle que : $1\text{ To} = 10^{12}\text{ octets}$ \quad ; \quad $1\text{ Go} = 10^9\text{ octets}$ \quad ; \quad $1\text{ Mo} = 10^6\text{ octets}$.

\begin{enumerate}
    \item Exprimer la capacité totale du disque dur en octets sous forme d'écriture scientifique : \pointille[4cm]
    \item Un fichier vidéo pèse $4\text{ Go}$. Combien de fichiers de ce type peut-on stocker sur ce disque dur ?
    \item Une connexion fibre télécharge à un débit de $100\text{ Mo/s}$ ($10^8\text{ octets/s}$).\\
    Calculer le temps nécessaire (en secondes) pour télécharger un jeu vidéo de $50\text{ Go}$.
\end{enumerate}
\lignesreponse[5]
\end{exercice}

% ------------------------------------------------------------------------------
% 7. TP TICE / CALCULATRICE
% ------------------------------------------------------------------------------
\section{TP TICE : Maîtriser les Puissances sur Calculatrice}

\begin{tpinfo}[Calculatrice Casio fx-92+ \& TI-Collège Plus]{Puissances et Notation Scientifique}
\begin{enumerate}
    \item \textbf{Saisir une puissance} :
    \begin{itemize}
        \item Pour calculer $3^5$ : taper \texttt{3} $\rightarrow$ touche puissance ($x^{\blacksquare}$ ou \texttt{\textasciicircum}) $\rightarrow$ \texttt{5} $\rightarrow$ \texttt{EXE}.
        \item Noter le résultat : $3^5 = \pointille[3cm]$
    \end{itemize}
    \item \textbf{Saisir une puissance de 10} :
    \begin{itemize}
        \item Utiliser la touche \texttt{$\times 10^x$} (ou \texttt{EE}) en bas du clavier.
        \item Pour entrer $4{,}2 \times 10^6$ : taper \texttt{4.2} $\rightarrow$ touche \texttt{$\times 10^x$} $\rightarrow$ \texttt{6} $\rightarrow$ \texttt{EXE}.
    \end{itemize}
    \item \textbf{Mode Scientifique (SCI)} :
    \begin{itemize}
        \item Configurer la calculatrice pour afficher directement en notation scientifique (Menu \texttt{CONFIG} $\rightarrow$ \texttt{Format affichage} $\rightarrow$ \texttt{Sci}).
        \item Calculer à la machine : $\dfrac{4{,}8 \times 10^7 \times 5 \times 10^{-2}}{2 \times 10^3} = \pointille[4cm]$
    \end{itemize}
\end{enumerate}
\end{tpinfo}

% ------------------------------------------------------------------------------
% 8. VERS LE BREVET (DNB PRO)
% ------------------------------------------------------------------------------
\section{Vers le Brevet (DNB Pro)}

\begin{sujetdnb}[DNB Pro Métropole][10 pts]{Vitesse de la Lumière \& Exploration Spatiale}
Dans le vide, la lumière et les ondes radio se propagent à la vitesse constante de :
\[ c = 3 \times 10^8\text{ m/s} = 300\,000\text{ km/s} \]
On rappelle la relation entre la vitesse $v$, la distance $d$ et la durée $t$ :
\[ v = \frac{d}{t} \iff d = v \times t \iff t = \frac{d}{v} \]

\begin{center}
\begin{tikzpicture}
    \node[draw=colDNB, fill=colDNBBg, rounded corners=6pt, line width=1pt, inner sep=8pt, text width=\dimexpr\linewidth-28pt\relax] {
        \small\sffamily
        \textbf{Données astronomiques :}
        \begin{itemize}[itemsep=1pt]
            \item Distance Terre -- Soleil : $d_1 = 1{,}5 \times 10^8\text{ km} = 1{,}5 \times 10^{11}\text{ m}$.
            \item Distance Terre -- Mars (au plus près) : $d_2 = 7{,}5 \times 10^7\text{ km} = 7{,}5 \times 10^{10}\text{ m}$.
        \end{itemize}
    };
\end{tikzpicture}
\end{center}

\begin{enumerate}
    \item \capSApproprier\ Exprimer la vitesse de la lumière $c$ en kilomètres par heure ($\text{km/h}$) sous forme d'écriture scientifique.
    \item \capRealiser\ Calculer le temps $t_1$ (en secondes) mis par la lumière du Soleil pour parvenir jusqu'à la Terre.
    \item \capRealiser\ Convertir ce temps $t_1$ en minutes et secondes.
    \item \capAnalyser\ Un robot explorateur situé sur Mars envoie un message radio vers la Terre.\\
    Combien de temps (en secondes, puis en minutes) s'écoule-t-il entre l'envoi du signal depuis Mars et sa réception par la station terrestre sur Terre ?
\end{enumerate}

\vspace{2mm}
\noindent\textbf{Rédigez votre démarche ci-dessous :}
\lignesreponse[10]
\end{sujetdnb}
